% УВОД, без номер. Ориентировъчен обем: 1 до 2 стр.
% Съдържа резюме на извършеното: цел, задачи, обхват, структура на изложението.
\section*{УВОД}
\label{sec:introduction}
\addcontentsline{toc}{section}{УВОД}

Сензорите за дълбочина и стереоскопичните камери извеждат данни, които не приличат на
обикновено растерно изображение. Резултатът е \term{облак от точки}: неподредено
множество от координати в тримерното пространство, чиито съседства не се четат от
индекса на пиксела, а трябва да бъдат намерени. Почти всяка стъпка от обработката,
от премахването на шума през оценяването на нормалите и извличането на признаци до
съвместяването на два припокриващи се записа, се свежда до един и същ въпрос: кои
точки лежат близо до дадена точка. Затова цената на целия конвейер се определя не
толкова от аритметиката на отделния алгоритъм, колкото от това колко бързо се отговаря
на този въпрос.

\textbf{Целта} на курсовия проект е да се проектира и реализира конвейер за обработка
на облаци от точки и растерни изображения, в който търсенето на съседи е изнесено в
отделен, ускорен компонент върху пространствен индекс, и всяка стъпка от конвейера е
измерима самостоятелно. Реализацията е на C++20 и се изгражда без нито една външна
зависимост. Ускорението идва от подредба на данните по компонента и от пакетно устроен
вътрешен цикъл, който компилаторът превежда във векторни инструкции; че това наистина се
случва, се проверява от отчета на компилатора, а не се предполага.

За постигането на целта са формулирани следните \textbf{задачи}:

\begin{enumerate}[topsep=0pt,itemsep=0pt,leftmargin=*]
\item преглед на апаратната и информационната структура на система за компютърно
      зрение и на класическите методи за предварителна обработка, филтриране,
      откриване на контури, извличане на признаци и сегментация;
\item анализ на възможните решения за пространствен индекс, за стратегия на
      векторизация и за организация на конвейера, и обосновка на направения избор;
\item проектиране на конвейер със самостоятелно измерими стъпки и с ясна граница
      между преносимата логика и векторизираната част;
\item програмна реализация на стъпките: четене и запис на облаци и изображения,
      прореждане с вокселна решетка, пространствен индекс със заявки за най-близки
      съседи и по радиус в два взаимозаменяеми варианта на вътрешния цикъл,
      съвместяване на облаци по схемата на итеративната най-близка точка, и растерна
      верига от изглаждане, градиент, потискане на немаксимуми и сегментация по
      свързани компоненти;
\item изграждане на методика за измерване, която сравнява всяка стъпка по точност и
      по време спрямо независим еталон;
\item провеждане на измерванията и анализ на получените резултати.
\end{enumerate}

\textbf{Обхватът} на проекта е геометричната и фотометричната обработка: всичко от
входните данни до съвместените облаци и сегментираните области. Извън него остават още
оценяването на дълбочина по стереодвойка и извличането на описания на околност: и двете
са потребители на заявката за съседи, чиято цена работата измерва пряко, и не участват в
измерваната верига. Извън обхвата
съзнателно остават разпознаването по обучени модели и невронните мрежи. Причината не е
липса на интерес, а несъвместима методика на оценяване: качеството на обучен
разпознавател зависи преди всичко от набора от данни и от процедурата на обучение, а
предмет на измерване тук е цената на геометричната обработка. Двете не се измерват с
един и същ експеримент, а обем от един курсов проект не стига и за двете. Тези теми са
разгледани в литературния преглед и мястото им в конвейера е посочено, но не са
реализирани.

\textbf{Структура на изложението.} Раздел~\ref{sec:literature} разглежда предметната
област и методите, върху които стъпва работата. Раздел~\ref{sec:alternatives} излага
разгледаните алтернативи и обосновава избора между тях. Раздел~\ref{sec:design}
представя проектирането на конвейера, а Раздел~\ref{sec:implementation} неговата
програмна реализация. Раздел~\ref{sec:method} описва методиката на измерване, а
Раздел~\ref{sec:results} получените резултати. Заключението обобщава постигнатото и
очертава посоките за продължение.
